\scanpage{19}
% L1
$S$ préschéma

% L2–5
$\Lambda$ faisceau loc. libre de rang fini de $\underline{O}_S$-algèbres $\underline{O}_S$-augmentées, \deleted{tel que} d'idéal d'augmentation $I$, tel que $\operatorname{gr}^{I}(\Lambda)$ soit encore localement libre [donc $I^n=0$ pour $n$ \unsure{grand|grd}]. Un tel $\Lambda$ sera dit ``\uline{faisceau de type infinitésimal}''

% L6
$T=\operatorname{Spec}(\Lambda)$, $S$-schéma fini et plat, muni d'une section $S\to T$

% L7–8
\[
\begin{aligned}
\Lambda_n&=\Lambda/I^{n+1},\quad T_n=\operatorname{Spec}(\Lambda_n)\quad\text{d'où}\\
T_0&\simeq S\subset T_1\subset T_2\ldots\subset T_n\subset\ldots T_N^{\deleted{\unsure{n}}}=T
\end{aligned}
\]

% L7–9, at left, diagram enclosed on three sides
\diagram{\begin{tikzcd}[ampersand replacement=\&,execute at end picture={\draw (current bounding box.north west) -- (current bounding box.north east) -- (current bounding box.south east) -- (current bounding box.south west);}]
\& X\arrow[d,"f"'] \& X_0\arrow[l]\arrow[d,"f_0"]\\
S \& T\arrow[l] \& S=T_0\arrow[l]
\end{tikzcd}}
% slanted
\margin{\unsure{donnons nous au} \unsure{préalable} $X_0$}

% L9
$X$ un $T$-\inserted{pré}schéma, \inserted{\uline{simple}} sur $T$, \quad $X_n=X\times_T T_n\deleted{\unsure{X}}$

% L10–11
On considère de plus
\[
\xi=X_0\mathbin{\deleted{\otimes}\times}_S T=X_0\otimes_{\underline{O}_S}\Lambda
\]
\margin{Donc $\xi_n=X_0\times_S T_n$ en particulier $\xi_0=X_0$}
% slanted
\margin{mettre les hypothèses pour \unsure{simplicité|simplifier}, \unsure{publier} en \unsure{Annales|tête} \unsure{pour} \unsure{l'énoncé|l'idée} de 2 \uline{\unsure{propositions|présentations}}}

% L12–13
Du point de vue espaces topologiques, on a
\[
|X|=|\xi|=|X_0|=|\xi_0|
\]

% L14–16
$\underline{O}_X$ est un faisceau \inserted{de $f_0^\ast(\underline{\Lambda})$-algèbres} sur $X_0$, augmenté vers $\underline{O}_{X_0}$,
\diagram{\begin{tikzcd}[ampersand replacement=\&]
\&\underline{O}_X\arrow[r] \& \underline{O}_{X_0}\\
\underline{O}_S\arrow[r] \& \underline{\Lambda}\arrow[u]\arrow[r,"\varepsilon"] \& \underline{O}_S\arrow[u]
\end{tikzcd}\quad avec commutativité}

% L17–18
de plus dans \deleted{\unsure{f|la}} filtration \inserted{$I$-adique} de $\underline{O}_X$, \deleted{\unsure{pour}} le gradué associé est loc. libre sur $\underline{O}_{X_0}$ ($X$ plat sur $T$).

% L19–21
\deleted{\unsure{Soit|Si}} Pour tt ouvert $U\subset X_0=|X|$, soit $\operatorname{Aut}^{X_0}_{T}(U)$ ou $\operatorname{Aut}^{X_0}_{\underline{\Lambda}}(X)$ le groupe des $T$-automorphismes de $U$ induisant l'identité sur $X_0\cap U=U\times_T T_0$.
% archive sheet number: 18

